\documentclass[11pt]{article}

\usepackage[margin=1in]{geometry}
\usepackage{amsmath,amssymb,bm}
\usepackage{booktabs}
\usepackage{graphicx}
\usepackage{microtype}
\usepackage{hyperref}
\usepackage{siunitx}
\usepackage{enumitem}
\usepackage{mathtools}

\newtheorem{theorem}{Theorem}
\newtheorem{corollary}{Corollary}
\newtheorem{proposition}{Proposition}
\newtheorem{remark}{Remark}

\newcommand{\Cfour}{\mathcal{C}_4}
\newcommand{\E}{\mathbb{E}}
\newcommand{\dd}{\mathrm{d}}

\title{Spectral-depth closure tests for falsifying photocarrier transport from Shockley--Ramo current}
\author{[Author]}
\date{Draft: August 10, 2026}

\begin{document}
\maketitle

\begin{abstract}
Photodetector frequency response is commonly interpreted by choosing a transport model and fitting material or circuit parameters to measured amplitude and phase. We ask a different question: can the detector be made to falsify simple transport hypotheses before a richer model is fitted? The required spatial coordinate can be supplied optically when wavelength moves the carrier-generation distribution through a monotonic absorber. Treating terminal-current formation explicitly with the Shockley--Ramo theorem, we show that a homogeneous one-carrier planar detector obeys an exact four-channel spatial closure,
\begin{equation*}
(J_2-J_1)^2=(J_1-J_0)(J_3-J_2),
\end{equation*}
for four equally spaced internal source coordinates. Spatial first differences isolate an internal propagation multiplier while rejecting wavelength-independent complex gain and additive offset. For uniform drift--diffusion with first-order recombination, four channels at DC plus the same four channels at one nonzero modulation frequency determine diffusion, drift, and recombination algebraically; every additional RF frequency introduces no new transport parameter and is therefore a null test. If the one-mode closure fails, six source coordinates test a two-mode model. The exact minor identity
\begin{equation*}
W_m=d_m d_{m+2}-d_{m+1}^2
=ab(q_1q_2)^m(q_1-q_2)^2
\end{equation*}
both supplies a rank-two closure and exposes the mode-resolution boundary. Finite scalar boundaries and conventional electron--hole transport then obey different RF root constraints. We derive leading optical-source, amplitude-calibration, internal-coordinate, and independent-noise corrections. In a deterministic high-Peclet limit the low-frequency four-channel phase closure measures a local combination of derivatives of inverse carrier velocity. A conditional graded-HgCdTe calculation predicts a gradient-sensitive closure phase of approximately \SI{-0.012}{\degree} at \SI{100}{MHz}, \SI{-0.059}{\degree} at \SI{500}{MHz}, and \SI{-0.110}{\degree} at \SI{1}{GHz}; an independent shooting implementation reproduces the stochastic calculation. The framework does not claim wavelength-dependent photodiode timing, Shockley--Ramo signal formation, Hankel model identification, or drift--diffusion inversion as new. Its purpose is to combine an internal spectral coordinate with observable-corrected spatial and RF closure relations into a deliberately falsification-driven transport measurement.
\end{abstract}

\section{Introduction}

A photodetector impulse or frequency response contains contributions from carrier generation, transport, recombination, boundaries, contacts, weighting-field geometry, and the external readout chain. A detailed model can therefore fit measured amplitude and phase without proving that its microscopic interpretation is unique. The objective here is not to eliminate modeling, but to change its order: begin with deliberately restrictive null hypotheses, derive exact closure relations that those hypotheses overdetermine, and introduce additional physical degrees of freedom only after a lower-order model fails.

Wavelength provides a natural internal coordinate in detectors where absorption depth changes strongly with photon energy. This physics is established. Optoelectronic chromatic dispersion uses wavelength-dependent absorption depth to create measurable photodiode RF phase shifts \cite{Glasser2021}; multi-frequency DC/RF feature extraction has also been used for single-photodiode computational spectroscopy \cite{Kassa2026}. Likewise, Hankel methods have been applied to time-sampled photodetector impulse responses for model identification \cite{Sun2025}. We do not claim these ingredients as new.

The narrower question is whether spectral channels can be calibrated as an \emph{internal spatial sequence}, converted into a Shockley--Ramo-aware spatial propagation sequence, and then subjected to model-order and RF-frequency closure tests. Three gedanken experiments organize the analysis:
\begin{enumerate}[label=(\roman*)]
    \item four internal source coordinates at one RF frequency test whether one homogeneous first-difference propagation mode is sufficient;
    \item the same coordinates at DC and another RF frequency test whether one real drift--diffusion--recombination generator survives a change in clock frequency;
    \item if one spatial mode fails, six coordinates test whether exactly two ordinary modes are sufficient and whether their RF roots match a finite boundary or conventional two-carrier explanation.
\end{enumerate}

The hierarchy is intentionally conservative. A failed four-channel law is not labeled anomalous transport: an ordinary electron--hole pair or a finite boundary is already sufficient to break it. A rank-two fit is not interpreted until the second-mode Hankel minor is statistically resolved. Only after these lower-order explanations are controlled do we interpret a residual in terms of spatially varying or richer transport.

HgCdTe is used only as a worked example. Wavelength/depth generation, graded-bandgap transport, and high-speed graded-HgCdTe response are established in the literature \cite{Sang2022}. Direct Shockley--Haynes measurements have independently characterized HgCdTe minority-carrier drift velocity, diffusion coefficient, and lifetime \cite{Rothman2010}. The material calculation below is a conditional prediction under explicit assumptions, not a fitted forecast for a named detector.

\section{Observable correction: first-passage flux versus terminal current}

\subsection{First-passage propagation}

Let one carrier begin a distance $d$ from a collecting boundary, and let $T_d$ be its successful first-passage time. Define
\begin{equation}
U(d,s)=\E[e^{-sT_d}].
\label{eq:Udef}
\end{equation}
For a homogeneous scalar first-passage process, spatial translation and the strong-Markov property imply
\begin{equation}
U(d_1+d_2,s)=U(d_1,s)U(d_2,s).
\end{equation}
Continuity in $d$ gives the spatial semigroup
\begin{equation}
\boxed{U(d,s)=e^{-\gamma(s)d}.}
\label{eq:semigroup}
\end{equation}
For uniform drift $w>0$, diffusion $D>0$, and uniform first-order killing/recombination rate $\kappa\ge 0$,
\begin{equation}
\boxed{D\gamma^2+w\gamma=\kappa+s.}
\label{eq:dispersion}
\end{equation}

Equation~\eqref{eq:semigroup} describes an arrival/collection-flux observable. It is not, by itself, the measured terminal current.

\subsection{Shockley--Ramo survival relation}

Let $p(x,t)$ be the density of carriers still present in a homogeneous one-dimensional region terminating at an absorbing collector $x=L$. Without killing,
\begin{equation}
\partial_t p=-\partial_x j,
\qquad
j=wp-D\partial_xp.
\label{eq:continuity}
\end{equation}
Define survival
\begin{equation}
S(t)=\int_{-\infty}^{L}p(x,t)\,\dd x.
\end{equation}
For a uniform planar weighting field $E_w$, Shockley--Ramo signal formation \cite{Shockley1938,Ramo1939,Dabrowski1989} gives
\begin{equation}
I(t)=qE_w\int_{-\infty}^{L}j(x,t)\,\dd x.
\end{equation}
If $p$ vanishes at the collector and remote upstream boundary, the diffusion boundary term cancels and
\begin{equation}
\boxed{I(t)=qE_w w S(t).}
\label{eq:ramo_survival}
\end{equation}
Since $\widetilde S(s)=[1-U(d,s)]/s$, the raw terminal-current transform is
\begin{equation}
J(d,s)=qE_w w\frac{1-U(d,s)}{s}.
\label{eq:Jnokill}
\end{equation}
Uniform Markov killing changes only the depth-independent prefactor. Using Eq.~\eqref{eq:semigroup},
\begin{equation}
\boxed{J(d,s)=C(s)\left[1-e^{-\gamma(s)d}\right].}
\label{eq:Jaffine}
\end{equation}

The distinction is not semantic. In deterministic uniform drift, the arrival transfer is $e^{-i\omega d/w}$, whereas the planar induced current is a rectangular pulse whose Fourier transform is proportional to $[1-e^{-i\omega d/w}]/(i\omega)$. A direct three-position geometric-mean law therefore applies to the ideal arrival field, not to generic terminal current. The measurable-current construction requires one additional spatial difference.

\section{Gedanken experiment I: four colors isolate one spatial mode}

Let four optical channels generate carriers at calibrated internal coordinates
\begin{equation}
d_m=d_0+mh,
\qquad m=0,1,2,3.
\label{eq:fourcoords}
\end{equation}
Equation~\eqref{eq:Jaffine} can be written
\begin{equation}
J_m=A+Bq^m,
\qquad q=e^{-\gamma h}.
\label{eq:affineseq}
\end{equation}
Define first differences $\Delta J_m=J_{m+1}-J_m$. Then
\begin{equation}
\Delta J_m=B(q-1)q^m.
\label{eq:firstdiff}
\end{equation}

\begin{theorem}[Four-color terminal-current closure]
For four equally spaced internal source coordinates in the homogeneous one-carrier planar model,
\begin{equation}
\boxed{(J_2-J_1)^2=(J_1-J_0)(J_3-J_2).}
\label{eq:fourclosure}
\end{equation}
\end{theorem}

Three source coordinates estimate
\begin{equation}
\boxed{q=\frac{J_2-J_1}{J_1-J_0},}
\label{eq:qrecover}
\end{equation}
while the fourth is a parameter-free falsification point. The spatial exponent follows from
\begin{equation}
\boxed{\gamma=-\frac{1}{h}\log q,}
\label{eq:grecover}
\end{equation}
with the logarithm branch tracked continuously in RF frequency or source coordinate.

\subsection{Exact nuisance invariances}

At one RF frequency, let
\begin{equation}
J_m^{\mathrm{meas}}=G(\omega)J_m+C(\omega),
\end{equation}
where $G$ and $C$ are common to the four spectral channels. First differences remove $C$ and multiply every $\Delta J_m$ by the same $G$, so Eq.~\eqref{eq:fourclosure} is unchanged.

Finite generation width is likewise not automatically a resolution failure. If each wavelength translates the same arbitrary generation shape $g$, averaging Eq.~\eqref{eq:Jaffine} changes only $B$; the sequence remains affine-exponential and Eq.~\eqref{eq:fourclosure} remains exact. The dangerous optical correction is wavelength-dependent \emph{shape evolution}, treated below.

Finally, if the true internal coordinate is related to a calibrated spectral coordinate by any affine map $z=a+b\mu$, equal spacing is preserved exactly. Absolute depth origin and common scale are unnecessary for the model-order null, although the scale is required to convert $\gamma$ to physical transport coefficients.

\section{Gedanken experiment II: DC plus one RF identifies; the next RF falsifies}

\subsection{No-recombination limit}

With $\kappa=0$, write $\gamma(i\omega)=a+ib$. Separating real and imaginary parts of Eq.~\eqref{eq:dispersion} yields
\begin{equation}
\boxed{D=\frac{\omega a}{b(a^2+b^2)},}
\label{eq:Dsimple}
\end{equation}
\begin{equation}
\boxed{w=\frac{\omega(b^2-a^2)}{b(a^2+b^2)}.}
\label{eq:wsimple}
\end{equation}
For the positive downstream branch,
\begin{equation}
\boxed{0<a<b.}
\end{equation}
One RF frequency therefore identifies $D$ and $w$ within the stated model; a second RF frequency introduces no new coefficient and must reproduce the same values.

\subsection{Complete inversion with recombination}

Use the four-channel sequence at DC to recover $g_0=\gamma(0)$ and at one RF to recover $g_\omega=\gamma(i\omega)$. The two dispersion equations are
\begin{align}
Dg_0^2+wg_0&=\kappa,\\
Dg_\omega^2+wg_\omega&=\kappa+i\omega.
\end{align}
Define
\begin{align}
A&=g_\omega^2-g_0^2,\\
B&=g_\omega-g_0,\\
\Delta&=\Re A\,\Im B-\Im A\,\Re B.
\end{align}
For $\Delta\neq0$,
\begin{equation}
\boxed{D=-\frac{\omega\Re B}{\Delta},}
\label{eq:Drecomb}
\end{equation}
\begin{equation}
\boxed{w=\frac{\omega\Re A}{\Delta},}
\label{eq:wrecomb}
\end{equation}
\begin{equation}
\boxed{\kappa=Dg_0^2+wg_0.}
\label{eq:krecomb}
\end{equation}

Thus four spectral channels at DC plus the same channels at one nonzero RF frequency determine the complete homogeneous three-parameter Markov model. Every additional RF frequency is overdetermined. The experiment can therefore be phrased as: \emph{DC plus one RF identifies the minimal model; the next RF tries to kill it.}

\section{Gedanken experiment III: six colors when one mode fails}

A failed four-color closure does not establish exotic transport. A finite boundary or an ordinary electron--hole pair can already generate two spatial modes.

Let five first differences from six source coordinates satisfy
\begin{equation}
d_m=a q_1^m+b q_2^m,
\qquad m=0,\ldots,4.
\label{eq:twomode}
\end{equation}
Define adjacent Hankel minors
\begin{equation}
W_m=d_md_{m+2}-d_{m+1}^2.
\end{equation}
Direct expansion gives
\begin{equation}
\boxed{W_m=ab(q_1q_2)^m(q_1-q_2)^2.}
\label{eq:Widentity}
\end{equation}
Hence
\begin{equation}
\boxed{W_1^2=W_0W_2,}
\label{eq:rank2closure}
\end{equation}
which is an exact six-color rank-two closure. Moreover,
\begin{equation}
\frac{W_{m+1}}{W_m}=q_1q_2,
\end{equation}
and the second-order recurrence yields $q_1+q_2$.

Equation~\eqref{eq:Widentity} also gives the structural identifiability boundary: evidence for two modes vanishes if either observable amplitude is zero and collapses quadratically as the two multipliers merge.

\subsection{Noise significance before root fitting}

For independent circular complex current noise with RMS $\sigma_J$,
\begin{equation}
\delta W_0
=-d_2\epsilon_0+(d_2+2d_1)\epsilon_1
-(d_0+2d_1)\epsilon_2+d_0\epsilon_3,
\end{equation}
so
\begin{equation}
\boxed{
\sigma_{W_0}^2=\sigma_J^2\left[
|d_2|^2+|d_2+2d_1|^2+|d_0+2d_1|^2+|d_0|^2
\right].}
\label{eq:Wnoise}
\end{equation}
Near equal current steps, $\sigma_{W_0}\simeq\sqrt{20}|d|\sigma_J$. For two comparable visible modes and $\eta=\sigma_J/|d|$,
\begin{equation}
Z_2\simeq\frac{|q_1-q_2|^2}{4\sqrt{20}\,\eta},
\end{equation}
so a $3\sigma$ design scale is
\begin{equation}
\boxed{|q_1-q_2|\gtrsim7.33\sqrt{\eta}.}
\label{eq:twores}
\end{equation}
A two-root fit should therefore not be attempted before the Hankel minor is statistically resolved.

\subsection{Finite scalar boundary}

For uniform scalar drift--diffusion--recombination with arbitrary linear finite-boundary conditions, the two homogeneous roots obey
\begin{equation}
D r^2+w r-(\kappa+i\omega)=0.
\end{equation}
Thus
\begin{equation}
\boxed{r_++r_-=-w/D,}
\label{eq:rootsum}
\end{equation}
which is real and RF-independent, and
\begin{equation}
\boxed{r_+r_-=-(\kappa+i\omega)/D.}
\label{eq:rootprod}
\end{equation}
The boundary amplitudes do not enter these constraints.

\subsection{Two conventional carrier species}

For independent homogeneous electron and hole contributions,
\begin{equation}
J(z,s)=C_0(s)+C_e(s)e^{+\gamma_e(s)z}+C_h(s)e^{-\gamma_h(s)z},
\end{equation}
where each propagation magnitude satisfies
\begin{equation}
D_c\gamma_c^2+w_c\gamma_c=\kappa_c+s.
\end{equation}
At DC the signed roots identify opposite collection directions in the generic recombining case. Tracking each root to one RF frequency and applying Eqs.~\eqref{eq:Drecomb}--\eqref{eq:krecomb} separately determines $(D_e,w_e,\kappa_e)$ and $(D_h,w_h,\kappa_h)$. Every later RF frequency overdetermines both triples.

The methodological point is that a failed one-mode model leads first to a model-order test, not immediately to a mechanism-specific fit.

\section{Controlled spatial inhomogeneity}

Once the one-mode observable and optical-source hypotheses have been controlled, a small closure residual can be connected to spatially varying transport.

Consider deterministic downstream transit with local slowness $q(z)=1/v(z)$ and uniform weighting field. A point-generated carrier has raw-current transform, up to a common factor,
\begin{equation}
J(z,s)=\int_z^L
\exp\left[-s\int_z^x q(u)\,\dd u\right]\dd x.
\label{eq:Jvarying}
\end{equation}
At low frequency,
\begin{equation}
J(z,s)=(L-z)-sA(z)+O(s^2),
\end{equation}
where $A(z)=\int_z^L(L-u)q(u)\,\dd u$. For four equally spaced source depths with spacing $h$ and midpoint $z_c$, define
\begin{equation}
\Cfour=2\ln\Delta J_1-\ln\Delta J_0-\ln\Delta J_2.
\end{equation}
A centered expansion gives
\begin{equation}
\boxed{
\Cfour=-s h^2\left[2q'(z_c)-(L-z_c)q''(z_c)\right]
+O(sh^4,s^2).}
\label{eq:gradienttheorem}
\end{equation}
For locally linear slowness,
\begin{equation}
\boxed{
\frac{\Im\Cfour}{\omega}=-2h^2q'(z_c)}
\label{eq:linearq}
\end{equation}
for the $e^{-i\omega t}$ convention.

Equation~\eqref{eq:linearq} is not an automatic interpretation of every nonzero closure residual. It is the positive prediction after lower-order source-shape, model-order, and high-Peclet assumptions have been controlled.

\section{Optical and calibration corrections}

\subsection{Source-shape evolution}

Let channel $m$ have generation coordinate $D_m=\mu_m+U_m$ with $\E[U_m]=0$, equally spaced means $\mu_m=\mu_0+mh$, and centered variances $v_m=\sigma_m^2$. Expanding the homogeneous raw-current closure gives
\begin{equation}
\boxed{
\Cfour_{\mathrm{opt}}
=\frac{\gamma}{2h}(v_3-3v_2+3v_1-v_0)+O(\gamma^2).}
\label{eq:opticalerror}
\end{equation}
Constant, linear, or quadratic variance evolution across the quartet therefore produces no first-order width contamination. Higher moments and higher powers of $\gamma$ remain real corrections.

\subsection{Relative amplitude calibration}

At sufficiently low RF, the ideal raw current is locally affine in source coordinate. If $\widetilde J_m=(1+\epsilon_m)J_m$,
\begin{equation}
\boxed{
\delta\Cfour=-\frac{\Delta^3(\epsilon_mJ_m)}{B},}
\label{eq:ampcal}
\end{equation}
where $B$ is the affine current step. Common fractional gain and linear fractional spectral drift cancel at first order; irregular or higher-curvature channel errors do not.

\subsection{Coordinate-map curvature}

For true coordinate $z=f(\mu)$,
\begin{equation}
\boxed{
\Cfour_{\mathrm{coord}}
=h^2\left[\gamma f''-(\ln f')''\right]_{\mu_c}+O(h^4).}
\label{eq:coorderror}
\end{equation}
Any affine map is therefore exact for the model-order null.

\subsection{Initial excess-energy invariance in an ideal graded gap}

Let $E_g(z)=E_{g0}-Gz$ and suppose local absorption depends only on photon excess energy $\nu=E_\gamma-E_g(z)$, so $\alpha=\alpha(\nu)$. Then the generation distribution expressed in $\nu$ is
\begin{equation}
\boxed{
p_\nu(\nu)
=\frac{\alpha(\nu)}{G}
\exp\left[-\frac{1}{G}\int_0^\nu\alpha(u)\,\dd u\right].}
\label{eq:energyinvariance}
\end{equation}
It is independent of photon energy in the untruncated ideal limit. Thus wavelength can translate generation depth while preserving the complete total excess-energy distribution. Real bandstructure and absorption only approximately satisfy these assumptions; the material example below quantifies the resulting mismatch.

\section{Independent-noise cost and spacing}

Linearizing the four-channel closure for independent circular complex errors gives
\begin{equation}
\delta\Cfour
=\frac{\epsilon_0}{d_0}
-\left(\frac1{d_0}+\frac2{d_1}\right)\epsilon_1
+\left(\frac2{d_1}+\frac1{d_2}\right)\epsilon_2
-\frac{\epsilon_3}{d_2}.
\end{equation}
In the equal-step limit,
\begin{equation}
\boxed{
\sigma_{\Cfour}\simeq\sqrt{20}\frac{\sigma_J}{|d|}.}
\label{eq:fournoise}
\end{equation}
The stencil is $(1,-3,3,-1)$: the same high-order cancellation that rejects smooth systematics amplifies uncorrelated noise.

If the dominant smooth systematic behaves as $Ah^2$ while statistical closure noise scales as $B/h$, the approximate mean-square error
\begin{equation}
\mathrm{MSE}(h)=A^2h^4+B^2/h^2
\end{equation}
has optimum
\begin{equation}
\boxed{h_*=\left(\frac{B}{\sqrt2A}\right)^{1/3}.}
\label{eq:hopt}
\end{equation}
Under white averaging, $\sigma_J\propto t^{-1/2}$ and therefore $h_*\propto t^{-1/6}$.

\section{Conditional graded-HgCdTe prediction}

\subsection{Optical coordinate}

As a worked example, consider a \SI{7.6}{\micro\meter} HgCdTe absorber at \SI{300}{K} with a linear composition profile from $x=0.55$ to $x=0.32$. The bandgap is evaluated with the Hansen--Schmit--Casselman relation \cite{Hansen1982}; above-gap absorption uses the Moazzami et al. model \cite{Moazzami2005}.

Four wavelengths are selected so that the conditional mean generation depths are
\begin{equation}
2.5,\ 3.0,\ 3.5,\ 4.0\ \si{\micro\meter}.
\end{equation}
The corresponding modeled wavelengths are approximately
\begin{equation}
2.134651,\ 2.215042,\ 2.301173,\ 2.393907\ \si{\micro\meter},
\end{equation}
and all four modeled absorbed fractions exceed $0.9993$. The generation widths are approximately \SI{0.79}{\micro\meter}, so the calculation is not a point-source stress.

The generation-weighted total photon excess energy above the local bandgap is also nearly matched. Its mean varies from approximately \SI{52.35}{meV} to \SI{52.48}{meV} and its standard deviation from approximately \SI{33.24}{meV} to \SI{32.69}{meV}, consistent with the invariance mechanism of Eq.~\eqref{eq:energyinvariance} while leaving electron--hole partition and thermalization memory as possible higher-order corrections.

\subsection{Transport stress}

Use a mobility sensitivity scale of \SI{9000}{\centi\meter^2\per\volt\per\second}, Einstein diffusion at \SI{300}{K}, a quasi-neutral composition-gap force scale, an \SI{8}{kV\per\centi\meter} velocity-saturation sensitivity scale, and the reduced density-of-states gradient correction developed in the model. These are explicit sensitivity coordinates, not claimed room-temperature constants for a particular detector.

The resulting downstream drift varies from approximately $3.76\times10^4$ to $3.21\times10^4\ \mathrm{m/s}$. Finite diffusion is retained through the backward equation for the expected discounted raw induced current,
\begin{equation}
\boxed{
D J''(z)+v(z)J'(z)-[\kappa+s]J(z)=-v(z).}
\label{eq:hgbvp}
\end{equation}
The collector condition is $J(L)=0$. To avoid the reflecting-boundary confound found in an earlier analysis, the entrance is matched mathematically to a bounded semi-infinite homogeneous continuation rather than modeled as a reflecting surface.

A homogeneous reference uses the same four optical kernels, diffusion coefficient, and recombination rate, but replaces the graded velocity by its path-harmonic mean. The difference between the graded and homogeneous four-channel closures is called the gradient-sensitive excess.

\subsection{Prediction and numerical cross-check}

For $\kappa=0$ the stochastic calculation gives
\begin{table}[h]
\centering
\caption{Conditional four-color HgCdTe closure prediction.}
\begin{tabular}{rrrr}
\toprule
RF & graded phase & same-optics homogeneous phase & excess \\
\midrule
\SI{100}{MHz} & \SI{-0.00952}{\degree} & \SI{+0.00246}{\degree} & \textbf{\SI{-0.01198}{\degree}} \\
\SI{500}{MHz} & \SI{-0.04643}{\degree} & \SI{+0.01230}{\degree} & \textbf{\SI{-0.05873}{\degree}} \\
\SI{1}{GHz} & \SI{-0.08572}{\degree} & \SI{+0.02470}{\degree} & \textbf{\SI{-0.11041}{\degree}} \\
\bottomrule
\end{tabular}
\label{tab:hgclosure}
\end{table}

The deterministic point-source low-frequency theorem, Eq.~\eqref{eq:gradienttheorem}, predicts approximately \SI{-0.01254}{\degree} at \SI{100}{MHz} for the same velocity profile and spacing, close to the full finite-width diffusion result. Uniform recombination stresses do not erase the effect in the tested range: an illustrative \SI{10}{ns} lifetime gives approximately \SI{-0.01205}{\degree} excess at \SI{100}{MHz}, while \SI{1}{ns} gives approximately \SI{-0.01275}{\degree}.

The stochastic result was independently reproduced with an adaptive shooting construction using a forced solution plus homogeneous boundary sensitivity. The sparse finite-difference boundary-value implementation and the shooting implementation agree to approximately $10^{-6}$ degree or better at \SI{100}{MHz}, \SI{500}{MHz}, and \SI{1}{GHz}.

These values are consequences of the stated conditional model, not forecasts for an existing device.

\subsection{Measurement resource}

For independent equal complex current noise, the exact four-sample covariance gives the following approximate $3\sigma$ requirements, expressed as amplitude SNR of the mean spatial current step:
\begin{table}[h]
\centering
\caption{Conditional $3\sigma$ measurement resource for the HgCdTe quartet.}
\begin{tabular}{rrr}
\toprule
RF & $\sigma_J/\langle|\Delta J|\rangle$ & amplitude SNR \\
\midrule
\SI{100}{MHz} & $1.56\times10^{-5}$ & \SI{96.1}{dB} \\
\SI{250}{MHz} & $3.88\times10^{-5}$ & \SI{88.2}{dB} \\
\SI{500}{MHz} & $7.68\times10^{-5}$ & \SI{82.3}{dB} \\
\SI{1}{GHz} & $1.47\times10^{-4}$ & \SI{76.7}{dB} \\
\bottomrule
\end{tabular}
\end{table}
The cleanest low-frequency asymptotic regime is therefore statistically demanding. Higher RF strengthens the closure signal in this explicit model, but also increases exposure to parasitic phase, additional spatial modes, and nonlocal dynamics.

\section{Discussion}

The main purpose of the closure hierarchy is to constrain interpretation order. A failed one-mode test should not immediately be interpreted as nonlocal transport. First ask whether a statistically resolved second spatial mode exists. If it does, test whether its roots obey a finite-boundary or conventional two-carrier law. Only when lower-order mechanisms fail should a richer transport state be introduced.

The framework sits next to several established bodies of work rather than replacing them. Shockley--Ramo signal formation is classical \cite{Shockley1938,Ramo1939,Dabrowski1989}. Wavelength-dependent absorption depth and photodiode RF phase are established through optoelectronic chromatic dispersion \cite{Glasser2021} and have already been used in multi-frequency computational spectroscopy \cite{Kassa2026}. Hankel model-order identification has been applied to time-sampled photodetector impulse responses \cite{Sun2025}. Uniform first-passage drift--diffusion and the quadratic propagation roots are standard stochastic-transport mathematics.

The candidate contribution considered here is therefore the \emph{combination}: wavelength is used as a calibrated internal spatial coordinate; raw Shockley--Ramo current is spatially differenced to isolate propagation modes; minimal color count tests model order; and the recovered spatial roots are forced to satisfy one physical law across RF frequency. Whether this exact spectral-depth construction has appeared previously remains an open priority question. A negative literature search is not evidence of priority, and no priority language is used here.

The strongest limitations are also clear. The exact lowest-rung theorems assume one-dimensional transport, a uniform planar weighting field, controlled relative generated-carrier amplitude, and a source kernel that is either rigidly translated or sufficiently well modeled. Nonuniform weighting fields, depletion-region signal formation, spatially varying coefficients, boundaries, carrier coupling, trapping, space charge, and nonlocal/hot-carrier dynamics can add spatial modes or RF dispersion. The hierarchy is meant to make those departures visible, not to assume them away.

\section{Conclusion}

Wavelength-dependent generation depth can be used for more than fitting a photodiode phase-versus-wavelength curve. When several spectral channels are calibrated as internal source coordinates, their complex terminal currents form a spatial sequence that can be subjected to exact model-order and transport-law closure tests.

Correct treatment of the electrical observable is essential. A homogeneous first-passage field is exponential in propagation distance, but a planar Shockley--Ramo terminal current contains a depth-independent particular term. Spatial first differencing removes that term. Four source coordinates then provide an exact one-mode closure and recover the spatial propagation exponent. DC plus one RF frequency determine the complete homogeneous drift--diffusion--recombination model; the next RF introduces no new transport parameter and therefore becomes a falsification measurement.

When one mode fails, six coordinates test whether exactly two observable spatial modes suffice. The exact minor $W_m=ab(q_1q_2)^m(q_1-q_2)^2$ both supplies the rank-two closure and exposes the mode-resolution limit. Finite boundaries and conventional two-carrier transport then make different RF root predictions. A richer physical model is warranted only if these lower-order closures fail.

The program is deliberately modest in what it reconstructs and ambitious in what it demands of a model. It does not seek an unrestricted pointwise map of every internal transport coefficient. It seeks measurements for which simple models make too many predictions---across color, model order, and RF frequency. The purpose is not to make transport fitting unnecessary, but to make the detector itself participate in deciding when a transport model has failed.

\bibliographystyle{unsrt}
\bibliography{MANUSCRIPT_REFERENCES}

\end{document}
